\documentclass[12pt,a4paper]{article}
\usepackage[fontset=none]{ctex} 
\setCJKmainfont{SimSun}

\usepackage[margin=1in]{geometry}
\usepackage{amsmath,amssymb}
\usepackage{hyperref}
\usepackage{fancyhdr}

\hypersetup{colorlinks,linkcolor=blue,citecolor=blue,urlcolor=blue}
\setlength{\headheight}{15pt}

\pagestyle{fancy}
\fancyhf{}
\fancyhead[L]{二维浅水方程 EnKF/EnSRF 理论说明}
\fancyhead[C]{libc\_DA\_swe}
\fancyhead[R]{\thepage}

\title{二维浅水方程的EnKF/EnSRF数据同化}
\author{杨鹏亮}
\date{\today}

\begin{document}

\maketitle

\begin{abstract}
本文按照 \texttt{swe2d\_enkf\_ensrf/swe2d\_enkf\_ensrf.c} 的程序实现，解释该程序所使用的二维浅水方程集合卡尔曼滤波（Ensemble Kalman Filter, EnKF）与集合平方根滤波（Ensemble Square Root Filter, EnSRF）方法。重点不是介绍最一般的集合滤波理论，而是说明这份代码中状态变量、观测算子、样本协方差、卡尔曼增益以及两种分析更新分别对应什么数学对象，以及代码中做了哪些近似和简化。
\end{abstract}

\tableofcontents

\section{引言}

这份代码实现的是一个最基本的“双系统”数据同化实验：
\begin{enumerate}
    \item 用同一个浅水方程数值模式推进一个“真值场”；
    \item 从真值场中抽取带噪声观测；
    \item 从同一初始集合复制出两套集合，分别用于随机 EnKF 和确定性 EnSRF；
    \item 在每个同化时刻，用同一批观测分别修正这两套集合；
    \item 对比两种滤波器的分析均值、RMSE 和集合离散度。
\end{enumerate}

因此，程序中的核心思想可以概括为
\[
\text{共同预测} \;\to\; \text{共同生成观测} \;\to\; \text{EnKF/EnSRF 分别分析} \;\to\; \text{继续预测}.
\]

\section{前向模式：二维浅水方程}

\subsection{连续方程}

代码的物理变量是
\[
h(x,y,t), \qquad hu(x,y,t), \qquad hv(x,y,t),
\]
其中 $h$ 是水深，$u,v$ 是速度分量，$hu,hv$ 是两个方向上的动量。

对应的二维浅水方程守恒形式为
\[
\frac{\partial \mathbf{q}}{\partial t}
+ \frac{\partial \mathbf{f}(\mathbf{q})}{\partial x}
+ \frac{\partial \mathbf{g}(\mathbf{q})}{\partial y}
= 0,
\]
其中状态向量为
\[
\mathbf{q} =
\begin{pmatrix}
h \\ hu \\ hv
\end{pmatrix}.
\]

在代码中使用的通量是
\[
\mathbf{f}(\mathbf{q}) =
\begin{pmatrix}
hu \\
\dfrac{(hu)^2}{h} + \dfrac{1}{2}gh^2 \\
\dfrac{hu\,hv}{h}
\end{pmatrix},
\qquad
\mathbf{g}(\mathbf{q}) =
\begin{pmatrix}
hv \\
\dfrac{hu\,hv}{h} \\
\dfrac{(hv)^2}{h} + \dfrac{1}{2}gh^2
\end{pmatrix}.
\]
这正对应代码中的 \texttt{flux\_x()} 和 \texttt{flux\_y()}。

\subsection{空间离散与时间推进}

在空间上，程序使用的是局部 Lax--Friedrichs，也叫 Rusanov 通量。以 $x$ 方向为例，两侧状态分别为 $\mathbf{q}_L,\mathbf{q}_R$ 时，数值通量写成
\[
\widehat{\mathbf{f}}_{i+1/2,j}
= \frac{1}{2}\left[\mathbf{f}(\mathbf{q}_L)+\mathbf{f}(\mathbf{q}_R)\right]
- \frac{1}{2}s_{\max}\left(\mathbf{q}_R-\mathbf{q}_L\right),
\]
其中
\[
s_{\max} = \max\left(|u_L|+\sqrt{gh_L},\; |u_R|+\sqrt{gh_R}\right).
\]
$y$ 方向同理。

于是一个显式 Euler 更新可以写成
\[
\mathbf{q}_{i,j}^{\,n+1}
= \mathbf{q}_{i,j}^{\,n}
- \frac{\Delta t}{\Delta x}\left(
\widehat{\mathbf{f}}_{i+1/2,j}-\widehat{\mathbf{f}}_{i-1/2,j}
\right)
- \frac{\Delta t}{\Delta y}\left(
\widehat{\mathbf{g}}_{i,j+1/2}-\widehat{\mathbf{g}}_{i,j-1/2}
\right).
\]
代码中这一步由 \texttt{euler\_update()} 完成。

在时间上，程序使用两步 RK2：
\begin{align*}
\mathbf{q}^{(1)} &= \mathbf{q}^n + \Delta t\, L(\mathbf{q}^n), \\
\mathbf{q}^{(2)} &= \mathbf{q}^{(1)} + \Delta t\, L(\mathbf{q}^{(1)}), \\
\mathbf{q}^{n+1} &= \frac{1}{2}\mathbf{q}^{n} + \frac{1}{2}\mathbf{q}^{(2)},
\end{align*}
其中 $L(\mathbf{q})$ 表示离散通量散度算子。代码中这一步由 \texttt{step\_rk2()} 完成。

\subsection{CFL 时间步长}

时间步长由 CFL 条件自适应给出：
\[
\Delta t
= \mathrm{CFL}\cdot
\min\left(
\frac{\Delta x}{|u|+\sqrt{gh}},
\frac{\Delta y}{|v|+\sqrt{gh}}
\right).
\]
代码里常数取为
\[
\mathrm{CFL}=0.40.
\]

\section{同化问题的状态、观测与集合表示}

\subsection{状态向量}

虽然代码内部把变量分别存储为 \texttt{h}, \texttt{hu}, \texttt{hv} 三个数组，但从 EnKF 的角度看，每个集合成员都对应一个大的状态向量
\[
\mathbf{x} =
\begin{pmatrix}
h_1,\dots,h_N,\;
(hu)_1,\dots,(hu)_N,\;
(hv)_1,\dots,(hv)_N
\end{pmatrix}^{\mathsf T},
\]
其中 $N=nx \times ny$ 是网格点总数。

因此程序实际同化的是整个三分量状态，而不仅仅是水深。

\subsection{集合初始化}

真值场由一个高斯隆起初始化：
\[
h(x,y,0) = 1 + 0.2\exp\left(-\frac{(x-x_0)^2+(y-y_0)^2}{0.02}\right),
\qquad hu=0,\qquad hv=0.
\]

每个集合成员都从真值初值出发，再加入高斯随机扰动：
\begin{align*}
h^{(e)} &= h^{\mathrm{truth}} + \sigma_{h0}\,\xi_h, \\
(hu)^{(e)} &= (hu)^{\mathrm{truth}} + \sigma_{m0}\,\xi_{hu}, \\
(hv)^{(e)} &= (hv)^{\mathrm{truth}} + \sigma_{m0}\,\xi_{hv},
\end{align*}
其中 $\xi$ 是标准正态随机变量。代码中取
\[
\sigma_{h0}=0.02,\qquad \sigma_{m0}=0.005.
\]

这一步的意义是：EnKF 不是维护单个状态，而是维护一个随机样本集合，用集合离散地表示概率分布。

\section{观测模型}

\subsection{观测网络}

程序在区域内部构造 $m=16$ 个观测点，近似为一个规则稀疏网格。观测点索引由 \texttt{build\_observation\_network()} 给出。

\subsection{观测算子}

这份代码只观测水深 $h$，不直接观测动量 $hu,hv$。因此观测方程可写为
\[
\mathbf{y} = \mathcal{H}(\mathbf{x}) + \boldsymbol{\varepsilon},
\]
其中观测算子 $\mathcal{H}$ 的作用只是“从完整状态向量中抽取若干个网格点上的 $h$ 值”。

如果第 $p$ 个观测对应网格点 $k_p$，那么
\[
\bigl[\mathcal{H}(\mathbf{x})\bigr]_p = h_{k_p}.
\]

观测误差取独立同分布高斯噪声：
\[
\boldsymbol{\varepsilon}\sim \mathcal{N}(0,R),
\qquad
R = \sigma_{\mathrm{obs}}^2 I.
\]
代码中
\[
\sigma_{\mathrm{obs}} = 0.01.
\]

因此观测生成公式是
\[
y_p = h^{\mathrm{truth}}_{k_p} + \sigma_{\mathrm{obs}}\eta_p,
\qquad
\eta_p\sim \mathcal{N}(0,1).
\]
这一步由 \texttt{make\_observations()} 完成。

\section{共同预测步}

设第 $e$ 个集合成员在某个同化时刻之前的分析状态为 $\mathbf{x}_{k-1}^{a,(e)}$。经过模式积分后得到预测状态
\[
\mathbf{x}_{k}^{f,(e)} = \mathcal{M}\bigl(\mathbf{x}_{k-1}^{a,(e)}\bigr),
\]
其中 $\mathcal{M}$ 就是由浅水方程数值模式定义的非线性推进算子。

在这份代码里，每次同化窗口内都推进固定步数 \texttt{assim\_every=20}。因此程序层面写成
\[
\mathbf{x}_{k}^{f,(e)} = \mathcal{M}_{20}\bigl(\mathbf{x}_{k-1}^{a,(e)}\bigr).
\]

随后由预测集合计算集合均值
\[
\bar{\mathbf{x}}^f = \frac{1}{N_e}\sum_{e=1}^{N_e}\mathbf{x}^{f,(e)},
\]
这里 $N_e$ 是集合成员数，代码中取
\[
N_e = 20.
\]

\section{随机 EnKF 分析步}

\subsection{这些公式从哪里来：从 KF 到随机 EnKF}

前面已经说明代码如何构造样本协方差，但如果只写
\[
K = P_{xy}(P_{yy}+R)^{-1},
\qquad
\mathbf{x}^{a,(e)}=\mathbf{x}^{f,(e)}+K\mathbf{d}^{(e)},
\]
确实会让人看不出这些公式的来源。这里补上随机 EnKF 最核心的一条理论链。下面这一组公式对应 Asch, Bocquet, Nodet
\emph{Data Assimilation: Methods, Algorithms, and Applications}
（2016）第 6.3 节，特别是式 (6.3)--(6.11)。

先从线性高斯 Kalman 滤波的分析步写起。若观测算子在分析时刻写成线性形式
\[
\mathbf{y} = H\mathbf{x} + \boldsymbol{\varepsilon},
\qquad
\boldsymbol{\varepsilon}\sim\mathcal{N}(0,R),
\]
且预测误差协方差为 $P^f$，则标准 Kalman 分析均值更新为
\[
\mathbf{x}^a = \mathbf{x}^f + K(\mathbf{y}-H\mathbf{x}^f),
\]
其中
\[
K = P^f H^{\mathsf T}\left(H P^f H^{\mathsf T}+R\right)^{-1}.
\]
对应的分析协方差满足
\[
P^a = (I-KH)P^f(I-KH)^{\mathsf T} + K R K^{\mathsf T}
= (I-KH)P^f.
\]

随机 EnKF 的目标不是直接维护 $P^f$，而是用有限集合来模仿这一步分析。设预测集合为
\[
\left\{\mathbf{x}^{f,(e)}\right\}_{e=1}^{N_e},
\qquad
\bar{\mathbf{x}}^f=\frac{1}{N_e}\sum_{e=1}^{N_e}\mathbf{x}^{f,(e)}.
\]
定义归一化离差矩阵
\[
X^f
= \frac{1}{\sqrt{N_e-1}}
\begin{bmatrix}
\mathbf{x}^{f,(1)}-\bar{\mathbf{x}}^f &
\cdots &
\mathbf{x}^{f,(N_e)}-\bar{\mathbf{x}}^f
\end{bmatrix},
\]
则样本预测协方差可写成
\[
P^f \approx X^f (X^f)^{\mathsf T}.
\]
这就是书里把协方差写成异常矩阵乘积的关键一步：EnKF 不显式推进庞大的 $P^f$，而是只推进集合成员本身。

如果直接对每个成员套用
\[
\mathbf{x}^{a,(e)}=\mathbf{x}^{f,(e)}+K(\mathbf{y}-H\mathbf{x}^{f,(e)}),
\]
那么成员离差满足
\[
X^a = (I-KH)X^f,
\]
从而样本分析协方差变成
\[
P^a_{\mathrm{sample}}
= (I-KH)P^f(I-KH)^{\mathsf T}.
\]
这比真正的 Kalman 分析协方差少了
\[
K R K^{\mathsf T}
\]
这一项，也就是低估了分析离散度。书中指出，这正是简单“同一观测直接更新所有成员”会导致集合过度收缩的原因。

因此随机 EnKF 改为给每个成员加入一份独立观测扰动
\[
\mathbf{y}^{(e)}_{\mathrm{pert}}
= \mathbf{y} + \mathbf{u}^{(e)},
\qquad
\mathbf{u}^{(e)}\sim \mathcal{N}(0,R),
\]
并执行
\[
\mathbf{x}^{a,(e)}
= \mathbf{x}^{f,(e)}
+ K\left(\mathbf{y}^{(e)}_{\mathrm{pert}}-H\mathbf{x}^{f,(e)}\right).
\]
在这种写法下，分析离差除了 $(I-KH)X^f$ 以外，还多出由观测扰动诱导的随机项；对这些随机项取期望，就会把缺失的
\[
K R K^{\mathsf T}
\]
补回来，因此
\[
\mathbb{E}[P^a_{\mathrm{sample}}]
= (I-KH)P^f(I-KH)^{\mathsf T}+K R K^{\mathsf T}
= (I-KH)P^f.
\]
这就是“扰动观测”并不是拍脑袋加噪声，而是为了让有限集合分析协方差在期望上满足 Kalman 理论。

再往下，把观测离差也写成矩阵形式。令
\[
\bar{\mathbf{y}}^f = \frac{1}{N_e}\sum_{e=1}^{N_e}H\mathbf{x}^{f,(e)},
\qquad
\bar{\mathbf{u}} = \frac{1}{N_e}\sum_{e=1}^{N_e}\mathbf{u}^{(e)},
\]
并定义归一化“创新离差”矩阵
\[
Y^f
= \frac{1}{\sqrt{N_e-1}}
\begin{bmatrix}
H\mathbf{x}^{f,(1)}-\mathbf{u}^{(1)}-\bar{\mathbf{y}}^f+\bar{\mathbf{u}} &
\cdots &
H\mathbf{x}^{f,(N_e)}-\mathbf{u}^{(N_e)}-\bar{\mathbf{y}}^f+\bar{\mathbf{u}}
\end{bmatrix}.
\]
则书中给出的增益写法是
\[
K = X^f (Y^f)^{\mathsf T}\left(Y^f (Y^f)^{\mathsf T}\right)^{-1}.
\]
因为
\[
X^f (Y^f)^{\mathsf T}
\approx P^f H^{\mathsf T},
\qquad
Y^f (Y^f)^{\mathsf T}
\approx H P^f H^{\mathsf T} + R,
\]
所以上式就是 Kalman 增益
\[
K=P^f H^{\mathsf T}(H P^f H^{\mathsf T}+R)^{-1}
\]
的集合近似形式。

对一般非线性观测算子 $\mathcal{H}$，书中还给出一个常用近似：
\[
H\left(\mathbf{x}^{f,(e)}-\bar{\mathbf{x}}^f\right)
\approx
\mathcal{H}\left(\mathbf{x}^{f,(e)}\right)-\bar{\mathbf{y}}^f.
\]
这意味着不必显式构造观测算子的切线性矩阵 $H$，只要把每个集合成员直接送进完整观测算子 $\mathcal{H}$ 即可。这正是本程序的做法。

\subsection{观测空间样本}

对于每个预测成员，先映射到观测空间：
\[
\mathbf{y}^{f,(e)} = \mathcal{H}\bigl(\mathbf{x}^{f,(e)}\bigr).
\]
因为这里只观测 $h$，所以 $\mathbf{y}^{f,(e)}$ 就是第 $e$ 个成员在观测点上的水深值。

观测空间均值为
\[
\bar{\mathbf{y}}^f = \frac{1}{N_e}\sum_{e=1}^{N_e}\mathbf{y}^{f,(e)}.
\]

\subsection{样本协方差}

EnKF 的关键是用集合样本协方差近似真实协方差。程序中计算了两个矩阵：

\paragraph{状态与观测之间的互协方差}
\[
P_{xy}
= \frac{1}{N_e-1}
\sum_{e=1}^{N_e}
\left(\mathbf{x}^{f,(e)}-\bar{\mathbf{x}}^f\right)
\left(\mathbf{y}^{f,(e)}-\bar{\mathbf{y}}^f\right)^{\mathsf T}.
\]
这个矩阵回答的问题是：“某个状态分量的偏差，和某个观测位置水深的偏差是否相关？”

\paragraph{观测空间协方差}
\[
P_{yy}
= \frac{1}{N_e-1}
\sum_{e=1}^{N_e}
\left(\mathbf{y}^{f,(e)}-\bar{\mathbf{y}}^f\right)
\left(\mathbf{y}^{f,(e)}-\bar{\mathbf{y}}^f\right)^{\mathsf T}.
\]
这个矩阵描述预测集合在观测空间中的离散程度。

代码中的 \texttt{Pxy} 和 \texttt{Pyy} 就是这两个样本协方差。

如果把上一小节的异常矩阵记号专门改写到本程序的非线性观测算子 $\mathcal{H}$ 上，那么
\[
X^f
= \frac{1}{\sqrt{N_e-1}}
\begin{bmatrix}
\mathbf{x}^{f,(1)}-\bar{\mathbf{x}}^f &
\cdots &
\mathbf{x}^{f,(N_e)}-\bar{\mathbf{x}}^f
\end{bmatrix},
\]
\[
Y^f
= \frac{1}{\sqrt{N_e-1}}
\begin{bmatrix}
\mathbf{y}^{f,(1)}-\bar{\mathbf{y}}^f &
\cdots &
\mathbf{y}^{f,(N_e)}-\bar{\mathbf{y}}^f
\end{bmatrix},
\]
就有
\[
P_{xy}=X^f(Y^f)^{\mathsf T},
\qquad
P_{yy}=Y^f(Y^f)^{\mathsf T}.
\]
因此本代码显式双循环累加出来的 \texttt{Pxy}、\texttt{Pyy}，与书中异常矩阵公式是完全一致的，只是写法从矩阵乘法变成了逐项求和。

\subsection{卡尔曼增益}

有了样本协方差以后，就构造
\[
S = P_{yy} + R,
\]
以及卡尔曼增益
\[
K = P_{xy}S^{-1}.
\]

其含义很直观：如果某个状态分量与观测误差之间相关性强，而且观测误差协方差 $R$ 又不大，那么该状态分量就应该被观测更强烈地修正。

代码用 \texttt{invert\_matrix()} 对 $S$ 做直接求逆，然后显式计算 $K$。

把它和上一节的书本公式对照起来看，本程序实际上采用的是
\[
K = P_{xy}(P_{yy}+R)^{-1},
\]
这正是
\[
K = X^f (Y^f)^{\mathsf T}\left(Y^f (Y^f)^{\mathsf T}+R\right)^{-1}
\]
在“先显式形成样本协方差，再做矩阵求逆”的实现形式。

\subsection{扰动观测更新}

这份程序使用的是随机 EnKF（stochastic EnKF），也就是对每个成员使用带扰动的观测：
\[
\mathbf{y}^{(e)}_{\mathrm{pert}}
= \mathbf{y} + \boldsymbol{\varepsilon}^{(e)},
\qquad
\boldsymbol{\varepsilon}^{(e)} \sim \mathcal{N}(0,R).
\]

于是每个成员的创新向量为
\[
\mathbf{d}^{(e)}
= \mathbf{y}^{(e)}_{\mathrm{pert}} - \mathbf{y}^{f,(e)}.
\]

分析更新公式为
\[
\mathbf{x}^{a,(e)}
= \mathbf{x}^{f,(e)} + K\,\mathbf{d}^{(e)}.
\]

这份代码最值得注意的一点是：虽然只观测了 $h$，但由于 $K$ 是由完整状态与观测之间的互协方差构成的，所以分析步会同时更新
\[
h,\qquad hu,\qquad hv.
\]
这正是 EnKF 的核心能力：通过协方差把观测信息传播到未被直接观测的变量上。

从理论上看，这里的扰动观测还有第二层作用：它保证分析后的样本协方差在期望上满足 Kalman 分析协方差公式，而不是系统性偏小。也正因为这个原因，随机 EnKF 与后面的 EnSRF 形成了清晰对比：前者靠随机观测扰动补足 $K R K^{\mathsf T}$，后者则靠确定性的平方根离差变换达到同样的协方差收缩目标。

\section{确定性 EnKF（EnSRF）分析步}

\subsection{为什么需要平方根滤波}

随机 EnKF 的分析更新
\[
\mathbf{x}^{a,(e)}
= \mathbf{x}^{f,(e)} + K\,\mathbf{d}^{(e)}
\]
需要对每个成员加入一份独立的观测扰动。这样做的优点是公式直接，但缺点也很明显：即使理论上分析协方差应该收缩到某个确定值，有限集合下额外加入的观测噪声仍会引入采样误差。

EnSRF 的思想是：\textbf{不再对观测做随机扰动}，而是把分析更新拆成
\[
\text{均值更新} + \text{离差更新},
\]
并直接构造一个确定性的离差收缩公式，使分析样本协方差满足卡尔曼滤波所要求的收缩关系。

按照 Asch, Bocquet, Nodet 书中第 6.4 节的说法，这类方法通常统称为
\textbf{deterministic EnKF} 或 \textbf{ensemble square root Kalman filter}。也就是说，本程序里的 EnSRF 不是另一类完全不同的滤波器，而正是确定性 EnKF 的一种实现。

\subsection{确定性 EnKF 的基本公式}

先仍然使用预测集合均值与归一化离差矩阵
\[
\bar{\mathbf{x}}^f=\frac{1}{N_e}\sum_{e=1}^{N_e}\mathbf{x}^{f,(e)},
\qquad
X^f
= \frac{1}{\sqrt{N_e-1}}
\begin{bmatrix}
\mathbf{x}^{f,(1)}-\bar{\mathbf{x}}^f &
\cdots &
\mathbf{x}^{f,(N_e)}-\bar{\mathbf{x}}^f
\end{bmatrix}.
\]
于是样本预测协方差仍写成
\[
P^f \approx X^f (X^f)^{\mathsf T}.
\]

在线性观测情形，令
\[
Y^f = H X^f.
\]
对一般非线性观测算子 $\mathcal{H}$，则像前面随机 EnKF 一样，用观测离差矩阵近似：
\[
Y^f
= \frac{1}{\sqrt{N_e-1}}
\begin{bmatrix}
\mathcal{H}(\mathbf{x}^{f,(1)})-\bar{\mathbf{y}}^f &
\cdots &
\mathcal{H}(\mathbf{x}^{f,(N_e)})-\bar{\mathbf{y}}^f
\end{bmatrix}.
\]

确定性 EnKF 的第一步是先更新均值，而不是先对每个成员单独更新。均值分析仍然服从 Kalman 公式
\[
\bar{\mathbf{x}}^a
= \bar{\mathbf{x}}^f + K\left(\mathbf{y}-\mathcal{H}(\bar{\mathbf{x}}^f)\right),
\]
其中
\[
K = P^f H^{\mathsf T}(H P^f H^{\mathsf T}+R)^{-1}.
\]
把 $P^f=X^f(X^f)^{\mathsf T}$ 代入，并改写到集合子空间后，书中给出的均值系数向量满足
\[
\bar{\mathbf{x}}^a = \bar{\mathbf{x}}^f + X^f \mathbf{w}^a,
\]
\[
\mathbf{w}^a
= \left(I + (Y^f)^{\mathsf T} R^{-1} Y^f\right)^{-1}
(Y^f)^{\mathsf T} R^{-1}
\left(\mathbf{y}-\mathcal{H}(\bar{\mathbf{x}}^f)\right).
\]
这说明确定性 EnKF 与随机 EnKF 在\textbf{均值}上并没有本质分歧；关键区别在于后面如何生成分析离差。

\subsection{平方根离差变换}

确定性 EnKF 的核心要求是：构造一个确定性的分析离差矩阵 $X^a$，使它满足
\[
P^a = X^a (X^a)^{\mathsf T}
= (I-KH)P^f.
\]
书中给出的集合子空间平方根公式可写成
\[
X^a = X^f T_a,
\]
其中右乘变换矩阵 $T_a$ 满足
\[
T_a T_a^{\mathsf T}
= \left(I + (Y^f)^{\mathsf T} R^{-1} Y^f\right)^{-1}.
\]
一个标准选取是对称平方根
\[
T_a
= \left(I + (Y^f)^{\mathsf T} R^{-1} Y^f\right)^{-1/2}.
\]

于是分析集合成员可写成
\[
\mathbf{x}^{a,(e)}
= \bar{\mathbf{x}}^a + \sqrt{N_e-1}\, [X^a]_e,
\]
其中 $[X^a]_e$ 表示 $X^a$ 的第 $e$ 列。这个公式说明：\textbf{确定性 EnKF 的本质，就是先更新分析均值，再对预测离差施加一个平方根变换，使样本协方差恰好收缩到理论上的分析协方差。}

从这个角度看，随机 EnKF 和确定性 EnKF 的差别可以概括为：
\begin{enumerate}
    \item 随机 EnKF 通过“扰动观测”在统计意义上补足 $K R K^{\mathsf T}$；
    \item 确定性 EnKF 则直接构造 $X^a$，使 $X^a(X^a)^{\mathsf T}$ 精确满足分析协方差关系。
\end{enumerate}

\subsection{本程序使用的是串行标量观测 EnSRF}

上面是确定性 EnKF 的一般矩阵形式。本程序中的 \texttt{ensrf\_analysis()} 没有显式构造
\[
T_a=\left(I + (Y^f)^{\mathsf T} R^{-1} Y^f\right)^{-1/2},
\]
也不是一次性处理整个观测向量，而是按观测点
\[
p=1,2,\dots,m
\]
逐个顺序更新。也就是说，它实现的是\textbf{串行标量观测 EnSRF}，可看成确定性 EnKF 在“每次只吸收一个标量观测”时的状态空间版本。对第 $p$ 个观测，记其对应网格索引为 $k_p$，则该标量观测为
\[
y_p = h_{k_p}^{\mathrm{truth}} + \varepsilon_p,
\qquad
\varepsilon_p \sim \mathcal{N}(0,R), \qquad R=\sigma_{\mathrm{obs}}^2.
\]

在这个标量观测场景里，矩阵平方根问题退化成一个一维离差收缩问题，因此最后会得到代码中的标量系数
\[
\alpha = \frac{1}{1+\sqrt{R/(H P^f H^{\mathsf T}+R)}}.
\]

由于这里只观测水深，所以对第 $e$ 个预测成员有
\[
h x_p^{(e)} = h_{k_p}^{f,(e)}.
\]
这里把 $hx$ 记成“观测算子作用到状态后的结果”，因为在一般写法里它对应
\[
h x_p^{(e)} = H_p \mathbf{x}^{f,(e)}.
\]

\subsection{第 1 步：观测均值与观测离差}

对固定的第 $p$ 个观测，程序先计算预测观测均值
\[
\bar{y}_p^f
= \frac{1}{N_e}\sum_{e=1}^{N_e} h x_p^{(e)},
\]
并构造每个成员的观测离差
\[
h x_p^{\prime (e)}
= h x_p^{(e)} - \bar{y}_p^f.
\]

代码中的数组 \texttt{hxp[e]} 就是这里的
\[
h x_p^{\prime (e)}.
\]

\subsection{第 2 步：标量观测方差}

然后代码计算该观测方向上的预测样本方差
\[
H P^f H^{\mathsf T}
= \frac{1}{N_e-1}
\sum_{e=1}^{N_e}
\left(h x_p^{\prime (e)}\right)^2.
\]

在程序里这个量记为 \texttt{hpht}。再加上观测误差方差，就得到标量分母
\[
\text{denom} = H P^f H^{\mathsf T} + R.
\]

\subsection{第 3 步：均值更新}

第 $p$ 个观测的创新写成
\[
d_p = y_p - \bar{y}_p^f.
\]

对任意一个状态分量 $x_r$，程序先计算它与该观测的样本协方差
\[
\mathrm{cov}(x_r, h x_p)
= \frac{1}{N_e-1}
\sum_{e=1}^{N_e}
\left(x_r^{f,(e)}-\bar{x}_r^f\right)
h x_p^{\prime(e)}.
\]

于是标量卡尔曼增益为
\[
K_r = \frac{\mathrm{cov}(x_r, h x_p)}{H P^f H^{\mathsf T}+R}.
\]

均值更新就是普通卡尔曼滤波的标量形式：
\[
\bar{x}_r^a = \bar{x}_r^f + K_r d_p.
\]

在代码里，对每个网格点的三个物理分量分别对应
\[
K_h,\qquad K_{hu},\qquad K_{hv},
\]
因此虽然只观测了水深，三个分量的均值都会通过协方差被同时修正。

\subsection{第 4 步：离差更新}

EnSRF 的关键区别在于：它不再使用扰动观测，而是直接更新成员离差。程序中采用的缩放系数为
\[
\alpha = \frac{1}{1+\sqrt{R/(H P^f H^{\mathsf T}+R)}}.
\]

于是第 $e$ 个成员的离差更新写成
\[
x_r^{\prime a,(e)}
= x_r^{\prime f,(e)} - \alpha K_r h x_p^{\prime (e)}.
\]

把均值更新与离差更新合并回成员本身，就得到代码中的公式
\[
x_r^{a,(e)}
= x_r^{f,(e)} + K_r d_p - \alpha K_r h x_p^{\prime (e)}.
\]

这正对应 \texttt{ensrf\_analysis()} 里的更新结构：
\[
\texttt{ens[e].state += mean\_increment - anomaly\_factor * hxp[e]}.
\]

\subsection{为什么这个系数是合理的}

对于标量观测，分析协方差应满足
\[
P^a = P^f - K H P^f.
\]
EnSRF 的目标是找到一个确定性的离差变换，使样本离差矩阵的协方差恰好满足这一关系。若把上一小节的一般平方根公式
\[
X^a = X^f T_a
\]
限制到“单个标量观测顺序吸收”的场景，那么矩阵平方根就退化成沿观测方向的一维收缩；上面的
\[
\alpha = \frac{1}{1+\sqrt{R/(H P^f H^{\mathsf T}+R)}}
\]
就是这一退化后的标准结果。换句话说，它正是确定性 EnKF 平方根变换在本程序这种串行标量实现下的具体表达式。它控制离差收缩的强度：
\begin{itemize}
    \item 当观测非常精确、$R$ 很小时，$\alpha$ 较大，离差收缩更明显；
    \item 当观测噪声较大时，$\alpha$ 较小，分析对集合离差的收缩也更弱。
\end{itemize}

\subsection{本程序中的串行实现含义}

因为程序是按观测点顺序更新，所以一次完整的 EnSRF 分析相当于做
\[
\mathbf{x}^{f}
\xrightarrow{\;y_1\;}
\mathbf{x}^{(1)}
\xrightarrow{\;y_2\;}
\cdots
\xrightarrow{\;y_m\;}
\mathbf{x}^{a}.
\]

理论上，如果线性高斯条件完全成立，并且数值上不受舍入影响，那么串行标量更新与一次性矢量更新是等价的。实际有限集合中，两者在实现细节和采样误差上会略有差别，但串行 EnSRF 的优点是实现简单，且避免了对整个观测协方差矩阵做一次大规模求逆。

\subsection{与随机 EnKF 的对比}

两种滤波器在本代码中的共同部分是：
\begin{enumerate}
    \item 使用同一个预测模式；
    \item 使用同一批观测；
    \item 都通过集合样本协方差把 $h$ 观测传播到 $hu,hv$；
    \item 分析后都做 $h \ge HMIN$ 的物理约束修正。
\end{enumerate}

它们的主要不同点是：
\begin{enumerate}
    \item 随机 EnKF 用完整观测向量一次更新，并显式构造
    \[
    K = P_{xy}(P_{yy}+R)^{-1};
    \]
    \item EnSRF 对每个标量观测顺序更新，不使用观测扰动；
    \item 随机 EnKF 的分析噪声部分来自人工加入的观测扰动；EnSRF 则通过确定性的离差缩放来达到正确的协方差收缩。
\end{enumerate}

\section{主程序中的同化循环}

把 \texttt{main()} 中的同化循环翻译成数学步骤，可以写成：

\subsection*{Step 1: 推进真值}
\[
\mathbf{x}^{\mathrm{truth}}_k
= \mathcal{M}_{20}\left(\mathbf{x}^{\mathrm{truth}}_{k-1}\right).
\]

\subsection*{Step 2: 推进两套集合成员}
\[
\mathbf{x}^{f,(e)}_{k,\mathrm{EnKF}}
= \mathcal{M}\left(\mathbf{x}^{a,(e)}_{k-1,\mathrm{EnKF}}\right),
\qquad
\mathbf{x}^{f,(e)}_{k,\mathrm{EnSRF}}
= \mathcal{M}\left(\mathbf{x}^{a,(e)}_{k-1,\mathrm{EnSRF}}\right),
\qquad e=1,\dots,N_e.
\]

\subsection*{Step 3: 分别计算两套预测集合均值}
\[
\bar{\mathbf{x}}^f_{k,\mathrm{EnKF}}
= \frac{1}{N_e}\sum_{e=1}^{N_e}\mathbf{x}^{f,(e)}_{k,\mathrm{EnKF}},
\qquad
\bar{\mathbf{x}}^f_{k,\mathrm{EnSRF}}
= \frac{1}{N_e}\sum_{e=1}^{N_e}\mathbf{x}^{f,(e)}_{k,\mathrm{EnSRF}}.
\]

\subsection*{Step 4: 从真值生成观测}
\[
\mathbf{y}_k
= \mathcal{H}\left(\mathbf{x}^{\mathrm{truth}}_k\right)
+ \boldsymbol{\varepsilon}_k.
\]

\subsection*{Step 5: 对 EnKF 构造矩阵增益}
\[
P_{xy,k}^{\mathrm{EnKF}},\qquad
P_{yy,k}^{\mathrm{EnKF}},\qquad
S_k = P_{yy,k}^{\mathrm{EnKF}}+R,\qquad
K_k = P_{xy,k}^{\mathrm{EnKF}}S_k^{-1}.
\]

\subsection*{Step 6: 对 EnKF 更新每个成员}
\[
\mathbf{x}^{a,(e)}_{k,\mathrm{EnKF}}
= \mathbf{x}^{f,(e)}_{k,\mathrm{EnKF}}
+ K_k
\left(
\mathbf{y}^{(e)}_{k,\mathrm{pert}} - \mathcal{H}(\mathbf{x}^{f,(e)}_{k,\mathrm{EnKF}})
\right).
\]

\subsection*{Step 7: 对 EnSRF 逐个观测顺序更新}
\[
\mathbf{x}^{f}_{k,\mathrm{EnSRF}}
\xrightarrow{\;y_{k,1}\;}
\mathbf{x}^{(1)}_{k,\mathrm{EnSRF}}
\xrightarrow{\;y_{k,2}\;}
\cdots
\xrightarrow{\;y_{k,m}\;}
\mathbf{x}^{a}_{k,\mathrm{EnSRF}}.
\]

\subsection*{Step 8: 计算两套分析均值}
\[
\bar{\mathbf{x}}^a_{k,\mathrm{EnKF}}
= \frac{1}{N_e}\sum_{e=1}^{N_e}\mathbf{x}^{a,(e)}_{k,\mathrm{EnKF}},
\qquad
\bar{\mathbf{x}}^a_{k,\mathrm{EnSRF}}
= \frac{1}{N_e}\sum_{e=1}^{N_e}\mathbf{x}^{a,(e)}_{k,\mathrm{EnSRF}}.
\]

程序输出文件可概括为：
\begin{itemize}
    \item \texttt{comparison\_output/truth\_cycle***.csv}：对应时刻的真值场；
    \item \texttt{comparison\_output/stoch\_enkf\_mean\_cycle***.csv}：随机 EnKF 的分析均值；
    \item \texttt{comparison\_output/ensrf\_mean\_cycle***.csv}：EnSRF 的分析均值；
    \item \texttt{comparison\_output/metrics.csv}：两种方法每个循环的 RMSE、spread 和 innovation 统计量。
\end{itemize}

\section{实现特点与局限}

\subsection{只观测水深，但更新完整状态}

虽然观测中只有 $h$，但动量分量 $hu,hv$ 仍然会被更新。这不是额外假设，而是由样本互协方差 $P_{xy}$ 自动给出的。如果某个网格点的动量变化总是和观测位置的水深变化相关，那么分析步就会修正该动量分量。

\subsection{没有局地化（localization）}

代码直接使用全局样本协方差，没有做 Schur 乘积局地化。因此远距离虚假相关可能会出现，尤其是在集合规模较小而状态维数很大时。

\subsection{没有膨胀（inflation）}

程序中没有对集合离差进行乘法膨胀，也没有附加模型误差。因此随着同化循环推进，集合离散度有可能逐渐偏小，出现“集合塌缩”。

\subsection{同时比较随机滤波与确定性平方根滤波}

程序把同一初始集合复制成两份：一份走随机 EnKF，一份走确定性 EnSRF。这样做的好处是比较公平，因为两种方法面对的真值轨道、观测网络、观测噪声水平和初始样本是一致的。两者差异主要来自分析步本身：EnKF 使用扰动观测，EnSRF 使用确定性的离差收缩。

\subsection{时间推进上的一个实现细节}

代码让真值和每个集合成员都各自根据 CFL 自适应步长推进固定的 \texttt{20} 个时间步，而不是推进完全相同的物理时间长度。因此严格来说，同一“同化时刻”下不同成员累计的物理时间可能略有差别。这个差别在数值上也许不大，但从理论上讲，这是当前实现中的一个近似。

\section{诊断量：RMSE\_f 与 RMSE\_a}

程序在每个同化循环中输出
\[
\mathrm{RMSE}_f, \qquad \mathrm{RMSE}_a,
\]
但这里是分别对 EnKF 和 EnSRF 各算一组，表示预测均值和分析均值相对于真值的水深均方根误差：
\[
\mathrm{RMSE}(h)
= \sqrt{
\frac{1}{N}
\sum_{k=1}^{N}\left(h_k^{\mathrm{mean}}-h_k^{\mathrm{truth}}\right)^2
}.
\]

通常希望看到
\[
\mathrm{RMSE}_a < \mathrm{RMSE}_f,
\]
这说明观测确实帮助修正了预测误差。若比较两种方法，则通常进一步关心
\[
\mathrm{RMSE}_{a,\mathrm{EnKF}}
\quad\text{与}\quad
\mathrm{RMSE}_{a,\mathrm{EnSRF}}
\]
谁更小，以及哪一种方法的分析误差随循环推进更稳定。

\section{结语}

对这份代码而言，可以只记住下面五点：
\begin{enumerate}
    \item 前向模式是二维浅水方程的有限体积离散，状态变量是 $(h,hu,hv)$。
    \item 观测只包含少量网格点上的水深 $h$，并带有高斯噪声。
    \item 随机 EnKF 用样本协方差构造矩阵卡尔曼增益，并通过扰动观测更新每个成员。
    \item EnSRF 不加观测扰动，而是把同化拆成“均值更新 + 离差确定性收缩”。
    \item 即使没有直接观测 $hu,hv$，它们也会因为与 $h$ 的样本相关性而被两种滤波器一并更新。
\end{enumerate}

如果把程序看成一句话，那么它做的就是：用两组共享初始条件的浅水方程样本分别实现随机 EnKF 和确定性 EnSRF，并在相同观测条件下比较它们把观测信息传播回整个流场的效果。

\appendix

\section{函数级实现细节：\texttt{ensrf\_analysis()}}

到这里为止，主文已经把理论主线、程序流程和诊断量交代完整。下面两节只做函数级的逐段公式对照，适合在同时打开 C 代码时查阅。

这一节把 \texttt{ensrf\_analysis()} 按代码执行顺序翻译成数学公式。它和上一节的 EnSRF 理论是一一对应的，只不过这里更强调实现细节。

\subsection{第 1 步：按观测顺序串行处理}

函数输入的是一整个集合 $\{\mathbf{x}^{f,(e)}\}_{e=1}^{N_e}$、观测向量 $\mathbf{y}$ 和观测索引 \texttt{obs\_idx}。但代码并不一次性处理所有观测，而是
\[
p = 1,2,\dots,m
\]
逐个更新。对每个观测点 $p$，它先取出对应的网格索引
\[
k_p = \texttt{obs\_idx[p]}.
\]

这意味着该函数实际实现的是串行标量 EnSRF：
\[
\mathbf{x}^{f}
\xrightarrow{\;y_1\;}
\mathbf{x}^{(1)}
\xrightarrow{\;y_2\;}
\cdots
\xrightarrow{\;y_m\;}
\mathbf{x}^{a}.
\]

\subsection{第 2 步：计算该观测的均值和离差}

对固定的第 $p$ 个观测，代码先计算观测均值
\[
\bar{y}_p^f = \frac{1}{N_e}\sum_{e=1}^{N_e} h_{k_p}^{f,(e)}.
\]

随后构造观测离差
\[
h x_p^{\prime(e)} = h_{k_p}^{f,(e)} - \bar{y}_p^f.
\]

在代码里，这个离差被临时存进数组 \texttt{hxp[e]}，以便后面的协方差和离差更新重复使用。

\subsection{第 3 步：计算标量分母}

接下来代码累加
\[
H P^f H^{\mathsf T}
= \frac{1}{N_e-1}\sum_{e=1}^{N_e}\left(h x_p^{\prime(e)}\right)^2,
\]
然后加上观测误差方差 $R=\sigma_{\mathrm{obs}}^2$，得到
\[
\text{denom} = H P^f H^{\mathsf T} + R.
\]

如果这个分母非常小，代码会跳过该观测，避免除零或不稳定更新。

\subsection{第 4 步：对每个状态分量计算协方差}

对任意状态分量 $x_r$，程序先计算集合均值，再构造样本协方差
\[
\mathrm{cov}(x_r, h x_p)
= \frac{1}{N_e-1}
\sum_{e=1}^{N_e}
\left(x_r^{f,(e)}-\bar{x}_r^f\right) h x_p^{\prime(e)}.
\]

在实现上，这里对三类分量分别计算：
\[
x_r \in \{h_k,\; (hu)_k,\; (hv)_k\}.
\]
这对应代码中的三个累加量 \texttt{cov\_h}, \texttt{cov\_hu}, \texttt{cov\_hv}。

\subsection{第 5 步：构造标量卡尔曼增益}

利用上面的协方差和分母，代码计算
\[
K_r = \frac{\mathrm{cov}(x_r, h x_p)}{H P^f H^{\mathsf T}+R}.
\]

如果把所有状态分量拼起来，这一串标量增益就构成了该观测对应的增益向量。但在代码里它并不显式形成整列向量，而是直接保留三个分量的标量增益用于后续更新。

\subsection{第 6 步：更新均值}

该观测的创新为
\[
d_p = y_p - \bar{y}_p^f.
\]

于是均值更新写成
\[
\bar{x}_r^a = \bar{x}_r^f + K_r d_p.
\]

在程序中，更新结果并不是先单独存均值，而是直接写回每个集合成员的状态更新式中。

\subsection{第 7 步：更新成员离差}

EnSRF 和随机 EnKF 的核心区别在这里。代码采用缩放系数
\[
\alpha = \frac{1}{1+\sqrt{R/(H P^f H^{\mathsf T}+R)}},
\]
并对每个成员执行
\[
x_r^{a,(e)}
= x_r^{f,(e)} + K_r d_p - \alpha K_r h x_p^{\prime(e)}.
\]

这等价于“均值修正 + 离差收缩”：
\[
x_r^{\prime a,(e)}
= x_r^{\prime f,(e)} - \alpha K_r h x_p^{\prime(e)}.
\]

代码里对每个网格点都分别算出 $h$、$hu$、$hv$ 的三个增量，然后一次性写回当前集合成员。

\subsection{第 8 步：物理约束检查}

完成该观测的更新后，代码调用
\[
\texttt{apply\_physical\_floor()}
\]
检查每个集合成员的水深是否满足
\[
h \ge HMIN.
\]

如果某个网格点出现负水深、非有限值或者动量发散，就把该点重置为
\[
h = HMIN,\qquad hu=0,\qquad hv=0.
\]

这一步和 \texttt{enkf\_analysis()} 中的处理是相同的，本质上是数值稳定性保护，不是滤波公式本身的一部分。

\subsection{一句话总结 \texttt{ensrf\_analysis()}}

如果把整个函数压缩成一句数学描述，那么它做的是：
\[
\text{对每个观测 } y_p,\;
\text{先算 } K_r \text{ 和 } \alpha,\;
\text{再用 } x_r^{a,(e)} = x_r^{f,(e)} + K_r d_p - \alpha K_r h x_p^{\prime(e)}
\text{ 顺序更新整个集合。}
\]

\section{函数级实现细节：\texttt{enkf\_analysis()}}

这一节的目标是把 \texttt{enkf\_analysis()} 从上到下翻译成数学公式。读完以后，应该可以直接对照 C 代码理解每一个数组在做什么。

\subsection{第 1 步：把三个物理数组拼成一个状态向量}

在代码里，单个成员的状态存成三段：
\[
\texttt{h[0:N-1]},\qquad \texttt{hu[0:N-1]},\qquad \texttt{hv[0:N-1]}.
\]
为了写成统一的矩阵公式，程序实际上等价于把它们拼成
\[
\mathbf{x}^{(e)} =
\begin{pmatrix}
\mathbf{h}^{(e)} \\
\mathbf{hu}^{(e)} \\
\mathbf{hv}^{(e)}
\end{pmatrix}
\in \mathbb{R}^{3N}.
\]

代码中
\[
\texttt{nstate} = 3\times \texttt{ncell}
\]
就是这个总状态维数。

\subsection{第 2 步：计算预测集合均值 \texttt{xmean}}

代码先对所有成员求平均：
\[
\bar{\mathbf{x}}^f
= \frac{1}{N_e}\sum_{e=1}^{N_e}\mathbf{x}^{f,(e)}.
\]

在实现上，\texttt{xmean} 的前三段分别存储
\[
\bar{\mathbf{h}}^f,\qquad \overline{\mathbf{hu}}^f,\qquad \overline{\mathbf{hv}}^f.
\]

也就是说，代码片段
\[
\texttt{xmean[k]},\quad
\texttt{xmean[ncell+k]},\quad
\texttt{xmean[2*ncell+k]}
\]
分别对应
\[
\bar h_k^f,\qquad \overline{hu}_k^f,\qquad \overline{hv}_k^f.
\]

\subsection{第 3 步：把每个成员映射到观测空间，得到 \texttt{Y} 与 \texttt{ymean}}

对于每个成员 $e$ 和每个观测点 $p$，代码执行
\[
Y[e,p] = h^{f,(e)}_{k_p},
\]
其中 $k_p=\texttt{obs\_idx[p]}$。这就是
\[
\mathbf{y}^{f,(e)} = \mathcal{H}(\mathbf{x}^{f,(e)}).
\]

随后求观测空间均值
\[
\bar{\mathbf{y}}^f
= \frac{1}{N_e}\sum_{e=1}^{N_e}\mathbf{y}^{f,(e)}.
\]

因此：
\begin{itemize}
    \item \texttt{Y[e*m+p]} 是第 $e$ 个成员在第 $p$ 个观测点上的预测观测值；
    \item \texttt{ymean[p]} 是该观测分量的集合均值。
\end{itemize}

\subsection{第 4 步：构造离差并计算 \texttt{Pxy}}

接下来代码进入一个三重循环。对每个成员和每个观测分量，它先计算观测离差
\[
dy_p^{(e)} = y_p^{f,(e)} - \bar y_p^f.
\]
然后对每个状态分量计算状态离差，例如
\[
dx_{h,k}^{(e)} = h_k^{f,(e)} - \bar h_k^f,
\]
\[
dx_{hu,k}^{(e)} = (hu)_k^{f,(e)} - \overline{hu}_k^f,
\]
\[
dx_{hv,k}^{(e)} = (hv)_k^{f,(e)} - \overline{hv}_k^f.
\]

代码中把它们逐项累加为
\[
P_{xy}[k,p] \mathrel{+}= dx_{h,k}^{(e)}\,dy_p^{(e)},
\]
\[
P_{xy}[N+k,p] \mathrel{+}= dx_{hu,k}^{(e)}\,dy_p^{(e)},
\]
\[
P_{xy}[2N+k,p] \mathrel{+}= dx_{hv,k}^{(e)}\,dy_p^{(e)}.
\]

最后再乘上
\[
\frac{1}{N_e-1}
\]
就得到标准样本互协方差
\[
P_{xy}
= \frac{1}{N_e-1}
\sum_{e=1}^{N_e}
\left(\mathbf{x}^{f,(e)}-\bar{\mathbf{x}}^f\right)
\left(\mathbf{y}^{f,(e)}-\bar{\mathbf{y}}^f\right)^{\mathsf T}.
\]

这一步非常关键，因为它正是观测信息传播到未观测变量的桥梁。

\subsection{第 5 步：计算 \texttt{Pyy}}

代码中
\[
\texttt{dyp = Y[e*m+p] - ymean[p]},\qquad
\texttt{dyq = Y[e*m+q] - ymean[q]}
\]
对应的是观测离差向量的两个分量。因此
\[
P_{yy}[p,q] \mathrel{+}= dyp^{(e)}\,dyq^{(e)}
\]
累加后并除以 $N_e-1$，得到
\[
P_{yy}
= \frac{1}{N_e-1}
\sum_{e=1}^{N_e}
\left(\mathbf{y}^{f,(e)}-\bar{\mathbf{y}}^f\right)
\left(\mathbf{y}^{f,(e)}-\bar{\mathbf{y}}^f\right)^{\mathsf T}.
\]

\subsection{第 6 步：构造 \texttt{S = Pyy + R}}

代码先令
\[
S = P_{yy},
\]
然后只在对角线上加上观测误差方差：
\[
S_{pp} \leftarrow S_{pp} + \sigma_{\mathrm{obs}}^2.
\]

因此这一步在数学上就是
\[
S = P_{yy} + R,
\qquad R = \sigma_{\mathrm{obs}}^2 I.
\]

之所以只加对角项，是因为程序假设不同观测误差彼此独立。

\subsection{第 7 步：求逆并计算卡尔曼增益 \texttt{K}}

代码调用
\[
\texttt{invert\_matrix(S, Sinv, m)}
\]
得到
\[
S^{-1}.
\]

随后按矩阵乘法显式累加
\[
K[k,q]
= \sum_{p=1}^{m} P_{xy}[k,p]\,(S^{-1})_{pq}.
\]

这正是
\[
K = P_{xy}S^{-1}.
\]

这里需要注意维数：
\[
P_{xy}\in \mathbb{R}^{3N\times m},\qquad
S^{-1}\in \mathbb{R}^{m\times m},\qquad
K\in \mathbb{R}^{3N\times m}.
\]

\subsection{第 8 步：生成扰动观测并计算创新 \texttt{innov}}

对每个成员，代码生成
\[
y^{(e)}_{\mathrm{pert},p}
= y_p + \sigma_{\mathrm{obs}} \eta_p^{(e)},
\qquad \eta_p^{(e)}\sim \mathcal{N}(0,1),
\]
然后计算创新
\[
\texttt{innov[p]} = y^{(e)}_{\mathrm{pert},p} - Y[e,p].
\]

也就是说
\[
\mathbf{d}^{(e)}
= \mathbf{y}^{(e)}_{\mathrm{pert}} - \mathbf{y}^{f,(e)}.
\]

\subsection{第 9 步：用 \texttt{K} 更新每个成员}

代码对每个网格点 $k$ 分别计算三个增量
\[
dh_k = \sum_{p=1}^{m} K_{h}(k,p)\, d_p^{(e)},
\]
\[
d(hu)_k = \sum_{p=1}^{m} K_{hu}(k,p)\, d_p^{(e)},
\]
\[
d(hv)_k = \sum_{p=1}^{m} K_{hv}(k,p)\, d_p^{(e)},
\]
然后更新
\[
h_k^{a,(e)} = h_k^{f,(e)} + dh_k,
\]
\[
(hu)_k^{a,(e)} = (hu)_k^{f,(e)} + d(hu)_k,
\]
\[
(hv)_k^{a,(e)} = (hv)_k^{f,(e)} + d(hv)_k.
\]

把三段合在一起，就是标准 EnKF 更新公式
\[
\mathbf{x}^{a,(e)}
= \mathbf{x}^{f,(e)} + K\,\mathbf{d}^{(e)}.
\]

\subsection{第 10 步：物理约束修正}

更新后，代码还检查
\[
h < HMIN
\]
的情形；如果出现，就把该点重置为
\[
h=HMIN,\qquad hu=0,\qquad hv=0.
\]

这一步不属于标准 EnKF 理论，而是数值稳定性与物理可接受性约束。它的作用是防止分析更新把局部水深改成非物理的负值或极小值。

\subsection{一句话总结 \texttt{enkf\_analysis()}}

如果把整个函数压缩成一句数学描述，那么它做的是：
\[
\text{用预测集合估计 } P_{xy},P_{yy},
\text{ 再构造 } K=P_{xy}(P_{yy}+R)^{-1},
\text{ 最后对每个成员执行 } \mathbf{x}^a=\mathbf{x}^f+K\mathbf{d}.
\]

\end{document}
